Let t be the trickster count and k the knight count. A knave can answer yes only if next to a trickster, so at most $2t$ knaves answer yes. Adding all non-knaves gives at most $3t+k\le3(t+k)$ yes answers. Since there are 17 yes answers, at least six people are not knaves, leaving at most 33 knaves.
For attainment, use six tricksters, all answering yes, separated cyclically by groups of $1,2,2,2,2,24$ knaves. Exactly eleven knaves are adjacent to a trickster and answer yes, while the other twenty-two answer no. Thus there are $6+11=17$ yes answers and 22 no answers, attaining $\boxed{33}$ knaves.
\begin{asy}
texpreamble("\usepackage{amsmath,amssymb}");
size(240);
defaultpen(fontsize(9pt));
string[] a={"T","L","T","L","L","T","L","L","T","L","L","T","L","L","T","L","L","L","L","L","L","L","L","L","L","L","L","L","L","L","L","L","L","L","L","L","L","L","L"};pair[] p;for(int i=0;i<39;++i)p.push(4dir(90-360.0*i/39));for(int i=0;i<39;++i){draw(p[i]--p[(i+1)%39],gray);label(a[i],p[i],a[i]=="T"?red:blue);}
\end{asy}
